% SOURCE: results/tables/heuristic_evaluation.md — "Defects found and
% corrected" (C1-C5) and the Synthesis; results/tables/heuristic_matrix.csv.
% ORDERING RULE, non-negotiable: DEFECTS FIRST. The distribution is reported
% once, in passing, and never as a headline or a score.
% Framing: single-evaluator, author-conducted, formative. Not user testing,
% not validation. Okolo/Solano-Kamaiko cited as why the user study is needed.
% BUDGET: 1 p. Cut prose before cutting a defect.
%
% 2026-09-08 framing pass (Decision Log #47). One paragraph added, giving the
% two instruments their own citations — the section named WCAG 2.1 and
% Nielsen's heuristics and cited neither, which left a methods claim standing
% on nothing. refs.bib gains nielsen1990heuristic, nielsen1994enhancing and
% w3c2025wcag21; all three resolved in-session. The "not user testing and not
% validation" paragraph below is UNCHANGED and must stay that way: the new
% paragraph gives the method standing, it does not upgrade what the method
% proves, and it says so in its own last two sentences.
%
% 2026-09-15 consolidated edit pass (see STATE.md):
%   (a) NEW table tab:matrix (Table 4 since 2026-10-09, when the gate table,
%       now Table 3, was added to Section 4; references use \ref), the 16 x 3
%       grid from heuristic_matrix.csv,
%       symbols only, rationales left in the CSV. The abstract already named
%       this matrix as archived and the paper had never shown it.
%   (b) Four/five disambiguated. Both numbers were right and referred to
%       different things: four Fail RATINGS on the first pass, arising from
%       five distinct DEFECTS (C1-C5 in heuristic_evaluation.md). "Four" is
%       NOT changed to "five"; a clause now says which is which.
%   (c) One editorial tail cut. The single-evaluator / not-user-testing /
%       not-validation paragraph is UNCHANGED, as its own header requires.
\section{Formative Heuristic Evaluation}
\label{sec:eval}

We rated the rendered artifacts against WCAG 2.1 and Nielsen's usability
heuristics: 16 criteria across three tiers, 48 cells, each grounded in a
named detail of a file on disk. Contrast ratios are computed at run time from
the palette constants rather than typed in by hand.

Both instruments are inspection methods, each with a stated procedure and a
literature on what it does and does not catch. Heuristic evaluation was
introduced as a way for an inspector rather than a user to find usability
problems against a small set of stated principles, and the set we rate
against is the refined ten~\citep{nielsen1990heuristic,
nielsen1994enhancing}; conformance inspection against the WCAG 2.1 success
criteria is the corresponding instrument on the accessibility
side~\citep{w3c2025wcag21}. Neither substitutes for users. Both were built to
find defects early.

This is a single-evaluator, author-conducted, formative assessment. It is not
user testing and it is not validation. Nothing in it is evidence that the
tiers help anyone. What it did was find defects in our own artifact, which is
what we report.

\subsection{What it found}

The first pass ran against the original renderer and returned four Fail
ratings, three of them on the health worker's card, arising from five distinct
defects. All five were fixed and the tiers re-rendered; the ratings below and
in Table~\ref{tab:matrix} describe the corrected artifacts.

\textbf{The driver list was sign-blind.} Features were ranked by absolute
SHAP magnitude with no sign filter, so a reading that argued \emph{against}
the prediction could be printed under a heading announcing it as a flag. On
the false-positive case the card listed low blood sugar, whose contribution
to the predicted class was $-0.085$. The routine card was worse: it listed
the two features that were the model's reasons for saying \emph{low} under
that same heading, above a next step telling the reader to arrange a
check-up. This is the defect that could have put a wrong cue in front of the
person making a referral, and no model metric would have shown it.

\textbf{The two tiers banded independently}, so the confident high-risk case
rendered an amber card beside a red lamp --- one prediction with two
identities. The band is now computed once and passed to both.

\textbf{Header contrast was assumed rather than measured.} White header text
was hardcoded; on the amber header, which is the state for two of the four
rendered cases, it measured $2.12{:}1$. Ink is now chosen per band by measured
contrast, worst case $5.14{:}1$, and a self-check asserts it.

\textbf{Direction words had no deadband}, so a blood sugar of $7.7$ against a
median of $7.5$ was worded as ``raised''. The word now attaches only outside
$0.25$ IQR of the median.

\textbf{The mother's transcript carried no provenance}, while both other
tiers printed the disclaimer. It now records case, band and disclaimer under
an explicit written-not-spoken rule; the spoken line is unchanged, because
appending provenance to every message would bury the one action the tier
exists to deliver.

\subsection{What it could not settle}

Table~\ref{tab:matrix} gives the full matrix. Across the 48 cells the
distribution after those fixes is 29 pass, 12 partial, 1 fail and 6 deferred.
The count is a description of coverage and not a score, and the five defects
above carry more weight. The single failure is the clinician tier shipping as
an image of text. The six deferred cells are all keyboard operability and
assistive-technology semantics, which a static PNG cannot exercise and a
deployed application would.

\begin{table}[t]
\centering
\footnotesize
\caption{All 48 cells: 16 criteria rated against each tier, on the corrected
artifacts. $\bullet$ pass, $\circ$ partial, $\times$ fail, -- deferred. The
per-cell rationale for every rating is in
\texttt{results/tables/heuristic\_matrix.csv}. Read by row the criteria are
covered; read by column each tier fails where a single universal rendering
would have to compromise.}
\label{tab:matrix}
\begin{tabular}{lccc}
\toprule
Criterion & Clinician & ASHA & Mother \\
\midrule
\multicolumn{4}{l}{\emph{WCAG 2.1}} \\
1.1.1 Text alternatives for non-text content & $\circ$ & $\bullet$ & $\circ$ \\
1.4.1 Colour not the only cue                & $\bullet$ & $\bullet$ & $\circ$ \\
1.4.3 Contrast (minimum)                     & $\bullet$ & $\bullet$ & $\bullet$ \\
1.4.5 Images of text                         & $\times$ & $\circ$ & $\bullet$ \\
2.1 / 2.4 Keyboard, focus order, navigation  & -- & -- & -- \\
3.1.1 / 3.1.2 Language of page and of parts  & $\bullet$ & $\circ$ & $\bullet$ \\
3.1.5 Reading level / plain language         & $\circ$ & $\circ$ & $\bullet$ \\
3.2.4 Consistent identification              & $\bullet$ & $\bullet$ & $\bullet$ \\
4.1.2 Name, role, value                      & -- & -- & -- \\
\addlinespace
\multicolumn{4}{l}{\emph{Nielsen's heuristics}} \\
1. Visibility of system status               & $\bullet$ & $\bullet$ & $\bullet$ \\
2. Match between system and the real world   & $\bullet$ & $\bullet$ & $\bullet$ \\
4. Consistency and standards                 & $\bullet$ & $\bullet$ & $\bullet$ \\
5. Error prevention                          & $\bullet$ & $\bullet$ & $\bullet$ \\
6. Recognition rather than recall            & $\circ$ & $\bullet$ & $\circ$ \\
8. Aesthetic and minimalist design           & $\circ$ & $\bullet$ & $\bullet$ \\
10. Help and documentation                   & $\circ$ & $\bullet$ & $\circ$ \\
\bottomrule
\end{tabular}
\end{table}

Read by column, each tier fails where a single universal rendering would have
to compromise: the clinician view needs chart and English literacy, the
mother's view needs hearing, and neither gap closes inside its own tier
without destroying what that tier is for. Read by row, the criteria are
covered. The architecture is what conforms; no single output does.

None of this tells us whether a health worker or a mother understands what
she is handed. That question needs users, and the nearest prior work is not
encouraging about assuming the answer: studies with community health workers
in India report difficulty interpreting AI explanations even after they have
been simplified~\citep{okolo2024easy,solanokamaiko2024explorable}. Studies with
mothers, health workers and clinicians are the next step
(Section~\ref{sec:future}).
